\textsc{Bimonads and Hopf monads} are a vast generalisation of bialgebras and Hopf algebras.
They naturally arise in the study of (rigid) monoidal categories and topological quantum field theories,
see amongst others~\cite{Kerler2001, Moerdijk2002, Bruguieres2007, Bruguieres2011, Turaev2017}.
Recall that the definition of a bialgebra object necessarily requires
a braided monoidal category as a base,
in order to write down what it means for the multiplication to be a morphism of comonoids.
However, in general the category of endofunctors is not braided,
so the naïve notion of bialgebras does not generalise to the monadic setting and needs to be adjusted.
One possible way of overcoming this problem was introduced and studied by Moerdijk under the name
\emph{Hopf monad},~\cite{Moerdijk2002}.\note{%
  As remarked in~\cite{Moerdijk2002},
  this concept is strictly dual to that of monoidal comonads,
  which are studied in~\cite{Boardman1995}.%
}
Here, the structure morphisms of
an oplax monoidal functor
serve as a substitute of the comultiplication and counit.
There are other, sometimes non-equivalent, notions of Hopf monad,
see~\cite{Boardman1995,Mesablishvili2011}.
We follow~\cite{Moerdijk2002},
with a slight terminology change due to~\cite{Bruguieres2007,Bruguieres2011}.
This definition aims to generalise the paradigmatic example of the free–forgetful adjunction of a Hopf algebra.
Thus, equipping a monad with the prefixes ``bi-'' or ``Hopf''
refers to additional structure or properties put on its Eilenberg–Moore category.

More generally, a monadic interpretation of module categories was given by Aguiar and Chase under the name \emph{comodule monad},~\cite{Aguiar2012}.
A comodule monad over a bimonad generalises the notion of a comodule algebra over a bialgebra,
see \cref{ex:comodule-algebra} below.
The main result of this chapter extends the reconstruction results
of~\cite[Proposition~4.1.2]{Aguiar2012} and~\cite[Lemma~7.10]{Turaev2017}:

\begin{reptheorem}{thm:comodule-adjunctions-correspond-to-strong-comodule-functors}
  Let \(\cat{C}\) and \(\cat{D}\) be monoidal categories,
  and suppose that \(\cat{M}\) and \(\cat{N}\) are right \(\cat{C}\)\hyp{} and \(\cat{D}\)\hyp{}module categories, respectively.
  Let \(\stdadj\) be an oplax monoidal adjunction.
  Lifts of an adjunction \(\adj{G}{V}{\cat{M}}{\cat{N}}\) to a comodule adjunction are in bijection with lifts of \(V \from \cat{N} \to \cat{M}\) to a strong comodule functor.
\end{reptheorem}

From the proof of this result one immediately obtains
an analogue of Kelly's doctrinal adjunction result,~\cite{kelly74:doctr},
for \(\cat{C}\)\hyp{}module categories.

\begin{reptheorem}{porism:doctrinal-module-adjunctions}
  An adjunction \(\adj{G}{V}{\cat{M}}{\cat{N}}\) of \(\cat{C}\)\hyp{}module categories
  yields a bijection between oplax \(\cat{C}\)\hyp{}module structures on \(F\)
  and lax \(\cat{C}\)\hyp{}module structures on \(U\).
\end{reptheorem}

We then obtain a Tannaka--Krein reconstruction result for comodule monads
in the spirit of~\cite[Theorem~7.1]{Moerdijk2002} and~\cite[Corollary~3.13]{McCrudden2002}.

\begin{reptheorem}{thm:comodule-monad-reconstruction}
  Let \(B\) be a bimonad on the monoidal category \(\cat{C}\),
  and \(T\) a monad on a right \(\cat{C}\)\hyp{}module category \(\cat{M}\).
  Coactions of \(B\) on \(T\) are
  in bijection with
  right actions of \(\cat{C}^B\) on \(\cat{M}^T\), such that \(U^{T}\) is a strict comodule functor over \(U^{B}\).
\end{reptheorem}

We can subsequently apply these results in order to study the comparison functor of \cref{sec:comp-funs}.
Namely, this functor always inherits the strong module structure of the ``strong adjoint'';
\ie, the right adjoint for oplax \(\cat{C}\)\hyp{}module adjunctions, and the left one in the lax case.

\begin{reptheorem}{strongcomparisons}
  Let \(\cat{M}\) and \(\cat{N}\) be left \(\cat{C}\)\hyp{}module categories.
  \begin{itemize}[noitemsep]
    \item The comparison functor \(K^{VG}\)
    of an oplax \(\cat{C}\)\hyp{}module adjunction \(\adj{G}{\!{}V}{\cat{M}}{\cat{N}}\)
    is a strong \(\cat{C}\)\hyp{}module functor.
    \item
    The comparison functor \(K_{VG}\) of
    a lax \(\cat{C}\)\hyp{}module adjunction
    \(\adj{G}{\!{}V}{\cat{M}}{\cat{N}}\)
    is a strong \(\cat{C}\)\hyp{}module functor.
  \end{itemize}
\end{reptheorem}

\section{Bimonads}\label{sec:monads-bimonads}

\begin{definition}\label{def:bimonad}
  A \emph{bimonad} on a monoidal category \(\cat{C}\) consists of
  an oplax monoidal endofunctor \(B\) on \(\cat{C}\), as well as
  oplax monoidal natural transformations \(\mu \from B^2 \nt B,\ \eta\from \Id_{\cat{C}} \nt B\)
  such that \((B, \mu, \eta)\) is a monad on \(\cat{C}\).

  A \emph{morphism of bimonads} is
  a natural transformation \(f\from B \nt H\) between bimonads
  that is oplax monoidal as well as a morphism of monads.
\end{definition}

Let us make \cref{def:bimonad} explicit.
In order for a monad \((B, \mu, \eta)\) to be a bimonad, there has to exist
a natural transformation \(B_2 \colon B \circ \otimes \nt B \otimes B\)
and a morphism \(B_0 \colon B1 \to 1\), such that all diagrams in \cref{fig:bimonad-explicitly} commute.
Further,
what it means for \(\mu\) and \(\eta\) to be oplax monoidal transformations
in the diagrammatic calculus of \cref{sec:string-diagr-mono}
is displayed in \cref{fig:mu-eta-opmonoidal-nt}.

\begin{figure}[htbp]
  \[\mathscale{0.95}{\begin{mytikzcd}[ampersand replacement=\&]
        {B (x \otimes y \otimes z)} \& {Bx \otimes B(y \otimes z)} \& {B1 \otimes Bx} \& Bx \& {Bx \otimes B1} \\
        {B(x \otimes y) \otimes Bz} \& {Bx \otimes By \otimes Bz} \&\& Bx
        \arrow["{B_{2; x,y \otimes z}}", from=1-1, to=1-2]
        \arrow["{x \otimes B_{2; y,z}}", from=1-2, to=2-2]
        \arrow["{B_{2; x \otimes y, z}}"', from=1-1, to=2-1]
        \arrow["{B_{2; x,y} \otimes z}"', from=2-1, to=2-2]
        \arrow["{B_{2; x,1}}", from=1-4, to=1-5]
        \arrow["{B_{2; 1,x}}"', from=1-4, to=1-3]
        \arrow["{\id_x}"', from=1-4, to=2-4]
        \arrow["{B_0 \otimes Bx}"', from=1-3, to=2-4]
        \arrow["{Bx \otimes B_0}", from=1-5, to=2-4]
      \end{mytikzcd}}\]
  \[\begin{mytikzcd}[ampersand replacement=\&]
      \& {B^2 (x \otimes y)} \&\& {B^2 1} \\
      {B(Bx \otimes By)} \&\& {B(x \otimes y)} \& B1 \\
      {B^2x \otimes B^2y} \&\& {Bx \otimes By} \& 1
      \arrow["{\mu_{x \otimes y}}", from=1-2, to=2-3]
      \arrow["{B_{2; x,y}}", from=2-3, to=3-3]
      \arrow["{BB_{2; x,y}}"', from=1-2, to=2-1]
      \arrow["{B_{2; Bx,By}}"', from=2-1, to=3-1]
      \arrow["{\mu_x \otimes \mu_y}"', from=3-1, to=3-3]
      \arrow["{\mu_{1}}"', curve={height=6pt}, from=1-4, to=2-4]
      \arrow["{B_0}"', from=2-4, to=3-4]
      \arrow["{B B_0}", curve={height=-6pt}, from=1-4, to=2-4]
    \end{mytikzcd}\]
  \[\begin{mytikzcd}[ampersand replacement=\&]
      {x \otimes y} \& {B(x \otimes y)} \&\& 1 \& B1 \\
      \& {Bx \otimes By} \&\& 1
      \arrow["{\eta_{x \otimes y}}", from=1-1, to=1-2]
      \arrow["{B_{2; x,y}}", from=1-2, to=2-2]
      \arrow["{\eta_x \otimes \eta_y}"', from=1-1, to=2-2]
      \arrow["{\id_{1}}"', from=1-4, to=2-4]
      \arrow["{\eta_{1}}", from=1-4, to=1-5]
      \arrow["{B_0}", from=1-5, to=2-4]
    \end{mytikzcd}\]
  \scaption[-6cm]{%
    Explicit conditions for \(B\) to be a bimonad.%
    \label{fig:bimonad-explicitly}%
  }
\end{figure}
\begin{figure}[htbp]
  \centering
  \mathscale{0.95}{\tikzfig{mu-eta-opmonoidal-nt}}
  \scaption[.5cm]{%
    Conditions for the multiplication and unit of a bimonad to be oplax monoidal natural transformations.%
    \label{fig:mu-eta-opmonoidal-nt}%
  }
\end{figure}

\begin{remark}\label{rmk:bimonad-in-bicategorical-language}
  We can also express \cref{def:bimonad} in more bicategorical language.
  Then, a bimonad may be defined as a monad\note{%
    A monoid in the monoidal category \(\BiCat{OplMon^{opl}}(\cat{C},\cat{C})\).%
  }
  in the bicategory \(\BiCat{OplMon^{opl}}\) where
  objects are monoidal categories,
  1-cells are oplax monoidal functors,
  and 2-cells are oplax monoidal transformations.
\end{remark}

\begin{example}\label{ex:bialgebra-gives-rise-to-bimonad}
  Let \(A \in \cat{C}\) be an object in a monoidal category.
  In \cref{ex:monoid-iff-tensoring-is-monad}
  we saw that the functor \(A \otimes \blank\) is a monad on \(\cat{C}\)
  if and only if \(A\) is an algebra object in \(\cat{C}\).
  This relationship extends to bialgebras:
  let \(B \in \cat{C}\) be an algebra object,
  and suppose \(T \defeq B \otimes \blank\) is the associated monad on \(\cat{C}\).
  Then \(T\) is a bimonad if and only if \(B\) is a bialgebra.
  More precisely, starting with a bialgebra structure \((\Delta,\varepsilon)\) on \(B\), set
  \[
    T_2\from B \otimes \blank \xrightarrow{\ \Delta \otimes \blank\ } B \otimes B \otimes \blank
    \qquad\text{and}\qquad
    T_0\from B \otimes \blank \xrightarrow{\ \varepsilon \otimes \blank\ } \blank.
  \]
  This turns \(T\) into a bimonad.

  Conversely,
  assume \((T, T_2, T_0)\) is an oplax monoidal functor.
  One defines a bialgebra structure on \(B\) by
  \[
    \Delta \from B = B \otimes 1 \otimes 1 \xrightarrow{T_{2; 1, 1}} B \otimes 1 \otimes B \otimes 1 = B \otimes B, \qquad
    \varepsilon \from B = B \otimes 1 \xrightarrow{\;T_0\;} 1.
  \]
  The fact that these arrows satisfy all of the required identities follows from a straightforward calculation.
\end{example}

\begin{example}\label{ex:oplaxmonoidal-adjunction-gives-bimonad}
  The intricate interplay between monads and adjunctions
  of \cref{ex:interplay-monad-adjunction}
  transcends to monoidal categories and bimonads.
  Given an oplax monoidal adjunction \(\stdadj\),
  the monad \(UF \from \cat{C} \to \cat{C}\)
  is a bimonad,
  whose comultiplication is defined for every \(x,y \in \cat{C}\) as the composition
  \[
    UF(x \otimes y)
    \trightarrow{UF_{2;x,y}}
    U(F x \otimes F y)
    \trightarrow{U_{2; F x,F y}}
    UF x \otimes UF y.
  \]
  Its counit is given by \(UF 1 \trightarrow{U F_0} U1 \trightarrow{\;U_0\;} 1\).
\end{example}

The following statement is a special case of~\cite[Theorem~1.2]{kelly74:doctr}.

\begin{proposition}\label{prop:doctrinaladjunction}
  Given monoidal categories \(\cat{C}\) and \(\cat{D}\),
  as well as an adjunction \(\adj{F}{U}{\cat{C}}{\cat{D}}\),
  there is a bijection between oplax monoidal structures on \(F\)
  and lax monoidal structures on \(U\).
\end{proposition}

The proof of~\cite[Theorem~1.2]{kelly74:doctr} is given constructively in terms of mates.
Given an oplax monoidal structure \((F_2, F_0)\) on \(F\),
the lax monoidal natural transformation \(U_2\) of \(U\) is,
for all \(x,y \in \cat{D}\),
given by
\[
  Ux \otimes Uy
  \xrightarrow{\eta_{Ux \otimes Uy}} UF(Ux \otimes Uy)
  \xrightarrow{UF_{2;Ux,Uy}} U(FUx \otimes FUy)
  \xrightarrow{U(\varepsilon_x \otimes \varepsilon_{y})} U(x \otimes y),
\]
where \(\eta\) and \(\varepsilon\) are the unit and counit of the adjunction.
Constructing an oplax monoidal structure on \(F\) from a lax monoidal structure on \(U\) is similar.

\newpage% XXX Sort of manual orphan prevention, as the below single line
        % paragraph should be "attached" to the following result more
        % than to the above paragraph.
Let us now highlight an important special case of \cref{prop:doctrinaladjunction}.

\begin{lemma}[{\cite[Lemma~7.10]{Turaev2017}}]\label{lem: monoidal_adjunctions_vs_bimonads}
  Any adjunction between two monoidal categories
  is oplax monoidal if and only if the right adjoint is strong monoidal.
\end{lemma}

\begin{example}
  Let \(\Bbbk\) be a field, and \(B \in \kVect\) a bialgebra.
  Then, by \cref{ex:oplaxmonoidal-adjunction-gives-bimonad},
  the monad \(B \kotimes \blank\from \kVect \to \kVect\) is a bimonad,
  and hence the adjunction \(\adj{B \kotimes \blank}{U}{\kVect}{B\textnormal{-Mod}}\)
  is oplax monoidal, where \(U\from B\textnormal{-Mod} \to \kVect\)
  is the canonical forgetful functor that simply forgets the \(B\)-module structure.
  By~\cref{lem: monoidal_adjunctions_vs_bimonads}, we obtain that \(U\) is even strong monoidal.

  More generally,
  let \(B \from \cat{C} \to \cat{C}\) be
  the bimonad arising from an oplax monoidal adjunction \stdadj.
  As the forgetful functor \(U^B \from \cat{C}^B \to \cat{C}\) is strict monoidal,
  the Eilenberg–Moore adjunction \(F^B \adjoint U^B\) is oplax monoidal.
\end{example}

The following result is due to~\cite{kelly74:doctr},
see also~\cite[Theorem~2.6]{Bruguieres2007}.

\begin{lemma}\label{lem: comparison_functor_bimonad_monoidal}
  Let \(\stdadj\) be an oplax monoidal adjunction.
  Then the comparison functor \(K^{UF}\from \cat{D} \to \cat{C}^{UF}\)
  of the bimonad \(UF\) is strong monoidal,
  and \(U^{UF} K^{UF} = U\) and \(K^{UF} F = F^{UF}\)
  as strong respectively oplax monoidal functors.
\end{lemma}

The question to which extent the monoidal structure on \(\cat{C}^B\) is unique
was answered
by Moerdijk~\cite[Theorem~7.1]{Moerdijk2002}
and McCrudden~\cite[Corollary~3.13]{McCrudden2002}.
In particular,
this kind of Tannaka--Krein reconstruction
gives another justification for the name ``bimonad''.

\begin{proposition}\label{prop: Moerdijk_reconstruction}
  Let \((B, \mu, \eta)\) be a monad on a monoidal category \(\cat{C}\).
  There exists a one-to-one correspondence between
  bimonad structures on \(B\) and
  monoidal structures on \(\cat{C}^B\)
  such that the forgetful functor \(U^B\) is strict monoidal.
\end{proposition}
\begin{proof}[Sketch of proof]
  Given a bimonad \((B, \mu ,\eta, B_2, B_0)\from \cat{C} \to \cat{C}\),
  as well as two modules \((m, \nabla_{m})\) and \((n, \nabla_{n}) \in \cat{C}^B\),
  we set
  \[
    (m, \nabla_{m}) \otimes (n, \nabla_n)
    \defeq
    \big(m \otimes n,\,
    B(m \otimes n) \xrightarrow{B_{2; m, n}} Bm \otimes Bn \xrightarrow{\nabla_m\otimes \nabla_n} m \otimes n
    \big).
  \]
  Moreover, define \(\nabla_{1}\from B1 \xrightarrow{\ B_0\ } 1\).
  The coassociativity and counitality of the comultiplication of \(B\)
  imply that the above construction implements a monoidal structure on \(\cat{C}^B\),
  parallel to that on the modules over a bialgebra.

  Conversely, let \((B, \mu, \eta)\) be a monad
  on the monoidal category \(\cat{C}\).
  Suppose \(\cat{C}^{B}\) is monoidal
  such that the forgetful functor \(U^B\) is strict monoidal.
  Consider the free modules \((Bm, \mu_m)\) and \((Bn, \mu_n)\),
  and denote their tensor product by
  \[
    (Bm, \mu_m) \otimes (Bn, \mu_n) = (Bm \otimes Bn, \mu_m \cdot \mu_n\from B(Bm \otimes Bn) \to Bm \otimes Bn).
  \]
  The comultiplication of \(B\) is then given by
  \[
    B(m \otimes n) \xrightarrow{B(\eta_m \otimes \eta_n)} B(Bm \otimes Bn) \xrightarrow{\mu_m \cdot \mu_n} Bm \otimes Bn.
  \]
  For the counit, take the action \(B1 \to 1\) of the unit of \(\cat{C}^{B}\).
\end{proof}

\section{Hopf monads}\label{sec:hopf-monads}

\textsc{We can now incorporate rigidity} into the framework of \cref{sec:monads-bimonads}.

\begin{definition}[{\cite[Section~2.6]{Bruguieres2011}}]\label{fusioninclusion}
  Let \(H\) be a bimonad on a monoidal category \(\cat{C}\).
  The \emph{right fusion operator} \(H_{\mathsf{rf}}\) of \(H\) is the natural transformation
  \[
    {H_{\mathsf{rf};x,y}}\from H(Hx \otimes y)
    \xrightarrow{H_{2;Hx,y}} H^{2}x \otimes Hy
    \xrightarrow{\mu_{x} \otimes Hy} Hx \otimes Hy, \quad\text{ for } x,y \in \cat{C}.
  \]
  The bimonad \(H\) is called \emph{right Hopf} if its right fusion operator is invertible.
\end{definition}

Left fusion operators and left Hopf monads are defined dually,
and a bimonad is said to be a \emph{Hopf monad} if it is both left and right Hopf.

\begin{theorem}[{\cite[Theorem~3.6]{Bruguieres2011}}]
  Let \(\cat{C}\) be a right closed monoidal category and let \(H\) be a bimonad on \(\cat{C}\).
  Then \(H\) is a right Hopf monad if and only if \(\cat{C}^H\) is right closed
  and the forgetful functor \(U^H\from \cat{C}^H \to \cat{C}\) is right closed.
\end{theorem}

A Hopf monad \(H\from \cat{C}\to \cat{C}\) can be defined on any monoidal category.
If \(\cat{C}\) is rigid then one may use~\cite[Lemma~3.4 and Theorem~3.6]{Bruguieres2011}
to obtain a rigid—not just closed—monoidal structure on the category of \(H\)\hyp{}modules.
This is reflected by the existence of two natural transformations
\begin{equation}\label{eq: left_and_right_antipode}
  s^{\ell}\from H(\ld{H}) \nt \ld{H}
  \qquad\text{and}\qquad
  s^{r}\from H(\rd{H}) \nt \rd{H}, \qquad \text{ for all }x \in \cat{C},
\end{equation}
called the \emph{left} and \emph{right antipode} of \(H\).
In Example~2.4 of~\cite{Bruguieres2012} it is explained how these generalise the antipode of a Hopf algebra;
an analogous result to \cref{ex:bialgebra-gives-rise-to-bimonad} holds in the Hopf case.

\section{(Co)module monads}\label{sec:comodule-monads}

\textsc{Alongside the theory of comodule monads} of Aguiar and Chase,~\cite{Aguiar2012},
we will also develop the special case of (op)lax \(\cat{C}\)\hyp{}module monads in this section.
The latter will play an important role in \cref{sec:infinite-non-rigid-reconstruction,sec:hopf-trimodules}.

\begin{definition}[{\cite[Definition~3.3.1]{Aguiar2012}}]\label{def: comodule_functor}
  Let \((F, F_2, F_0)\from \cat{C} \to \cat{D}\) be an oplax monoidal functor,
  \(\cat{M}\) a right \(\cat{C}\)-module,
  and \(\cat{N}\) a right \(\cat{D}\)-module category.
  A \emph{(right) comodule functor over \(F\)}\note{%
    Alternatively, we will say that \(G\) is a \emph{(right) \(F\)\hyp{}comodule functor}.%
  }
  is a pair \((G, G_{\mathsf{a}})\) consisting of
  a functor \(G \from \cat{M} \to \cat{N}\)
  together with a natural transformation
  \[
    G_{\mathsf{a}; m,x} \from G(m \ract x) \to G m \ract F x, \qquad \text{for all \(x\in\cat{C}\) and \(m \in \cat{M}\)},
  \]
  called the \emph{coaction} of \(G\),
  which is \emph{coassociative} and \emph{counital},
  in the sense that the following diagrams commute for all \(x,y \in \cat{C}\) and \(m \in \cat{M}\):
  \[
    \begin{mytikzcd}[ampersand replacement=\&]
	{G((m \ract x) \ract y)} \& {G(m \ract (x \otimes y))} \& {G(m \ract 1)} \& Gm \\
	{G(m \ract x) \ract Fy} \& {Gm \ract F(x \otimes y)} \& {Gm \ract F1} \& {Gm \ract 1} \\
	{(Gm \ract Fx) \ract Fy} \& {Gm \ract (Fx \otimes Fy)}
	\arrow["\cong", from=1-1, to=1-2]
	\arrow["{G_{\mathsf{a}; m\ract x, y}}"', from=1-1, to=2-1]
	\arrow["{G_{\mathsf{a}; m, x \otimes y}}", from=1-2, to=2-2]
	\arrow["\cong", from=1-3, to=1-4]
	\arrow["{G_{\mathsf{a};m,1}}"', from=1-3, to=2-3]
	\arrow["{G_{\mathsf{a};m,x} \ract Fy}"', from=2-1, to=3-1]
	\arrow["{Gm \ract F_{2;x,y}}", from=2-2, to=3-2]
	\arrow["{Gm \ract F_0}"', from=2-3, to=2-4]
	\arrow["\cong"', from=2-4, to=1-4]
	\arrow["\cong"', from=3-1, to=3-2]
      \end{mytikzcd}
    \]

  A comodule functor is called \emph{strong} if its coaction is an isomorphism,
  and \emph{strict} if it is the identity.
\end{definition}

\begin{example}\label{ex:comodule-functors-as-module-functors}
  A right comodule functor over the identity functor \(\Id_{\cat{C}}\) of a monoidal category \(\cat{C}\) is
  nothing more than an oplax \(\cat{C}\)\hyp{}module functor in the sense of \cref{def:lax-module-functor},
  considering right instead of left \(\cat{C}\)\hyp{}module categories.
\end{example}

\begin{example}[{\cite[Section~6.1]{Aguiar2012}}]\label{ex:comodule-algebra}
  If \(B\) is a bialgebra over a commutative ring \(\Bbbk\),
  then \(B \otimes_{\Bbbk} \blank \from \Bbbk\text{-Mod} \to \Bbbk\text{-Mod}\) becomes an oplax monoidal functor.
  Let \(A\) be a right \(B\)\hyp{}comodule algebra
  with coaction \(\nu \from A \to A \otimes_{\Bbbk} B\).
  Recall that the subalgebra of \emph{\(B\)-coinvariants} of \(A\) is given by
  \[
    A^{\mathrm{co}B} \defeq \{\, a \in A \mid \nu(a) = a \otimes 1 \,\}.
  \]
  Suppose that \(C\) is a subalgebra of \(A^{\mathrm{co}B}\).
  Then \(\nu\) is a map of \(C\)\hyp{}\(C\)\hyp{}bimodules,
  which turns \(A \otimes_{C} \blank \from C\text{-Mod} \to C\text{-Mod}\) into a right comodule functor over \(B \otimes_{\Bbbk} \blank\).
  The action of \(\Bbbk\textnormal{-Mod}\) on \(C\text{-Mod}\) is given by tensoring over \(\Bbbk\).
\end{example}

Recall the constructions of the centre and twisted centre from \cref{sec:drinfeld-centre,subsec: twisted_centres}.
The canonical functor \(U^{(Z)} \from \cent{C} \to \cat{C}\)
that forgets the half braiding is strict monoidal.
Similarly,
given a strong monoidal functor \(L\),
the forgetful functor \(U^{(L)} \from \cent[L]{C} \to \cat{C}\)
is a strict comodule functor over \(U^{(Z)}\).

\begin{remark}\label{rmk:graphical-comodule-functors}
  The diagrammatic calculus of \cref{sec:string-diagr-mono} can be adapted to the setting of comodule functors.
  In addition to combining sheets using tensor products,
  we now additionally consider actions to do so as well.
  In order to keep track of which splitting occurred,
  functors between the module categories shall be coloured blue.
  For example, a coaction will be drawn as:
  \[
    \tikzfig{graphical-coaction}
  \]

  The coassociativity and counitality of \(G_{\mathsf{a}}\) are displayed in \cref{fig:coassoc-and-counital-comodule-functor};
  notice the analogous nature of the diagrams to \cref{fig:coassoc-and-counital-opmonoidal-functor}.
  \begin{figure}[htbp]
    \centering
    \tikzfig{coassoc-and-counital-comodule-functor}
    \scaption[-1cm]{%
      Coassociativity and counitality conditions of the coaction of a comodule functor.%
      \label{fig:coassoc-and-counital-comodule-functor}%
    }
  \end{figure}
\end{remark}

\begin{definition}\label{def: morphism-of-comodules}
  Suppose that \(C, G \from \cat{M} \to \cat{N}\) are comodule functors over \(B, F \from \cat{C} \to \cat{D}\).
  A \emph{comodule (natural) transformation} from \(C\) to \(G\) comprises
  a pair of natural transformations \(\phi \from C \nt G\) and \(\psi \from B \nt F\),
  such that the following diagram commutes
  for all \(x \in \cat{C}\) and \(m \in \cat{M}\):
  \begin{equation}\label[diagram]{eq:comodule-natural-transformation}
    \begin{mytikzcd}[ampersand replacement=\&]
      {C(m \ract x)} \& {G(m \ract x)} \\
      {Cm \ract Bx} \& {Gm \ract Fx}
      \arrow["{\phi_{m \ract x}}", from=1-1, to=1-2]
      \arrow["{G_{\mathsf{a}; m,x}}", from=1-2, to=2-2]
      \arrow["{C_{\mathsf{a}; m,x}}"', from=1-1, to=2-1]
      \arrow["{\phi_m \ract \psi_x}"', from=2-1, to=2-2]
    \end{mytikzcd}
  \end{equation}

  We call \((\phi, \psi)\) a \emph{morphism of comodule functors}
  if \(B=F\) and \(\psi=\id_B\).
\end{definition}

\begin{example}\label{ex:module-trafos-as-comodule-trafos}
  If \(\cat{C} = \cat{D}\) and \(B = F = \Id_{\cat{C}}\),
  then \(C\) and \(G\) are oplax \(\cat{C}\)\hyp{}module functors, see \cref{ex:comodule-functors-as-module-functors}.
  A comodule natural transformation from \(C\) to \(G\) is exactly a \(\cat{C}\)\hyp{}module transformation
  in the sense of \cref{def:module-transformation}.
\end{example}

Given a comodule natural transformation \((\phi \from G \nt C,\, \psi \from B \nt F)\),
the graphical version of \cref{eq:comodule-natural-transformation} is displayed in our next picture,
where the blue dot represents \(\phi\) and the black dot represents \(\psi\).
\[
  \tikzfig{comodule-transformation}
\]

The notion of a morphism of comodule functors might seem restrictive;
however, it is actually sufficient to only consider arrows of this form.

\begin{example}\label{ex:comod-nat-trafo-is-comod-morphism}
  Let the pair \(\phi \from C \nt G\) and \(\psi \from B \nt F\) be a comodule transformation.
  We can view \(\phi \from C \nt G\) as a morphism of comodule functors over \(F\) if we equip \(C\) with a new coaction.
  It is given by
  \[
    C(m \ract x) \xrightarrow{\;C_{\mathsf{a};m,x}\;} Cm \ract Bx \xrightarrow{Cm\ract\psi_x} Cm \ract Fx,
    \qquad\text{for all \(x \in \cat{C}\) and \(m \in \cat{M}\).}
  \]
  Thus, by suitably altering the involved coactions,
  comodule natural transformations and morphisms of comodule functors can be identified.
\end{example}

Let \(B\from \cat{C} \to \cat{C}\) be a bimonad and \(\cat{M}\) a module category over \(\cat{C}\).
The unit \(\eta \from \Id_{\cat{C}} \nt B\) implements a coaction on \(\Id_{\cat{M}} \from \cat{M} \nt \cat{M}\) via
\[
  \id_m \ract \eta_x \from \Id_{\cat{M}}(m \ract x) \to \Id_{\cat{M}}m \ract Bx, \qquad \text{ for all } x \in \cat{C}, m \in \cat{M}.
  \]
Using the multiplication \(\mu \from B^2 \nt B\), we can equip the composition \(CG\) of two comodule functors \(C, G \from \cat{M} \to \cat{M}\) with a comodule structure:
\[
  {(CG)}_{\mathsf{a}}
  \from CG(\blank \ract \bblank)
  \xrightarrow{CG_{\mathsf{a}}} C(G(\blank) \ract B(\bblank))
  \xrightarrow{C_{\mathsf{a}}} CG(\blank) \ract B^2(\bblank)
  \xrightarrow{\id \ract \mu}
  CG(\blank) \ract B(\bblank).
\]
Due to the associativity and unitality of the multiplication of \(B\),
in this way the category of comodule endofunctors on \(\cat{M}\) over \(B\) becomes monoidal.

\begin{definition}\label{def: comodule_monad}
  Consider a bimonad \(B \from \cat{C} \to \cat{C}\)
  and a module category \(\cat{M}\) over \(\cat{C}\).
  A \emph{comodule monad} over \(B\) on \(\cat{M}\) comprises
  a comodule endofunctor \((C,C_{\mathsf{a}}) \from \cat{M} \to \cat{M}\) together with
  comodule morphisms \(\mu \from C^2 \nt C\)  and \(\eta \from \Id_{\cat{M}} \nt C\)
  such that \((C,\mu,\eta)\) is a monad.

  A \emph{morphism of comodule monads} is
  a natural transformation of comodule functors \(f \from C \nt G\)
  that is also a morphism of monads.
\end{definition}

The conditions for the multiplication and unit of
a comodule monad \(C\) on \(\cat{M}\)
over a bimonad \(B\) on \(\cat{C}\)
to be morphisms of comodule functors are given in \cref{fig:comodule-monad}.
Notice how the conditions are analogous to \cref{fig:mu-eta-opmonoidal-nt}.
\begin{figure}[htbp]
  \centering
  \tikzfig{comodule-monad}\vspace{-0.3cm}% XXX
  \scaption{%
    Conditions for the multiplication and unit of a comodule monad to be morphisms of comodule functors.%
    \label{fig:comodule-monad}%
  }
\end{figure}

\begin{example}\label{ex:oplax-module-monad}
  An \emph{oplax \(\mathcal{C}\)\hyp{}module monad} is a comodule monad over the identity monad on \(\cat{C}\).
  Alternatively, it is a monoid in the category
  \(\mathsf{Oplax}\mathcal{C}\mathsf{Mod}(\cat{M},\cat{M})\)
  of oplax \(\cat{C}\)\hyp{}module endofunctors on a (left) \(\cat{C}\)\hyp{}module category \(\cat{M}\).

  Analogously, one can define the notion of a \emph{lax \(\cat{C}\)\hyp{}module comonad}.
\end{example}

\begin{example}\label{ex:comodule-monads-example}
  Consider the poset \((\real, \leq)\).
  It is a monoidal category with addition as tensor product and \(0 \in \real\) as unit.
  In~\cite{hasegawa-lemay2018:LinearHopf}, Hasegawa and Lemay defined a bimonad using the ceiling function.
  Let
  \[
    H \from \real\to \real, \qquad x \mapsto \lceil x \rceil \defeq \min \{\, n\in \mathbb{Z} \mid n \geq x \,\}.
  \]
  As \(\lceil x + y \rceil \leq \lceil x \rceil + \lceil y \rceil\) for all \(x,y \in \real\) and \(\lceil 0 \rceil = 0\), the functor \(H\) is oplax monoidal.
  The comultiplication and counit are given by the unique arrows
  \[
    H_0\from H0 = 0 \to 0 \qquad\text{and}\qquad H_{2; x,y} \from H(x+y) \to Hx+ Hy,
  \]
  for all \(x,y \in \real\).
  The idempotence of the ceiling function implies that the identity natural transformation defines a multiplication \(\mu \from H^2 \nt H\).
  Its unit corresponds to \({\{\, x \to Hx \,\}}_{x \in \real}\).

  Given a number \(z\) in the half-open interval \([0, 1)\), let
  \[
    C^{(z)} \from \real \to \real, \qquad x \mapsto \min\{\, z+ n \mid n \in \mathbb{Z}, z+n \geq x \,\}.
  \]
  In case \(z=0\), we have \(C^{(z)} = H\).
  Otherwise \(C^{(z)}0 = z > 0\) and \(C^{(z)}\) cannot be oplax monoidal.
  Nonetheless, \(C^{(z)}(x+y) \leq C^{(z)} x + \lceil y \rceil = C^{(z)} x + Hy\) holds,
  and the unique natural arrow
  \[
    C^{(z)}_{\mathsf{a};x,y}\from C^{(z)}(x+y) \to C^{(z)} x + Hy, \qquad\qquad \text{for all } x,y\in \real,
  \]
  defines a coaction of \(H\) on \(C\).
  Again, we have that \({(C^{(z)})}^2 = C^{(z)}\) is idempotent
  and \(x \leq C^{(z)}x\), for all \(x \in \real\).
  Thus, it is a comodule monad over \(H\).
\end{example}

\begin{example}\label{ex:copower-as-comodule-monad}
  Let \((\cat{V}, \otimes, 1)\) be a closed symmetric monoidal category.
  A \(\cat{V}\)\hyp{}category \(\cat{C}\)
  is said to be \emph{copowered}
  over \(\cat{V}\),
  see~\cite[Section~3.7]{kelly05:basic},
  if there exists a functor \(\blank \cdot \bblank \from \cat{C} \times \cat{V} \to \cat{C}\),
  such that for all \(c \in \cat{C}\) we have
  \[
    \adj{c \cdot \blank}{\cat{C}(c, \blank)}{\cat{V}}{\cat{C}}.
  \]
  One obtains \(c\cdot (v \otimes w) \cong (c\cdot v) \cdot w\) by the Yoneda lemma:
  \begin{align*}
    \cat{C}\big(c \cdot (v \otimes w), x\big)
    & \cong \cat{V}\big(v \otimes w, \cat{C}(c, x)\big)
      \cong \cat{V}\big(w, \cat{V}(v, \cat{C}(c, x))\big) \\
    & \cong \cat{V}\big(w ,\cat{C}((c \cdot v), x)\big)
      \cong \cat{C}\big((c \cdot v)\cdot w, x\big).
  \end{align*}

  In fact, this turns \(\cat{C}\) into a \(\cat{V}\)\hyp{}module category,
  and therefore the identity monad on \(\cat{C}\) into a comodule monad over \(\Id_{\cat{V}}\) with trivial coaction.

  Suppose that \(B \defeq UF\) is a bimonad on \(\cat{V}\).
  The unit \(\eta \from \Id_{\cat{V}} \nt B\) is a morphism of bimonads and we may extend the coaction of \(\Id_{\cat{C}}\) to
  \[
    \delta\from \blank \cdot \bblank \xrightarrow{\; \blank \,\cdot\, \eta\;} \blank \cdot B\bblank,
  \]
  turning \(B\) into a \(\cat{C}\)-module monad.
\end{example}

The following proposition connects
the notion of a \(\cat{C}\)\hyp{}module monad
to that of an algebra object in \(\cat{C}\) in the sense of \cref{sec:algebra-and-module-objects}.

\begin{proposition}\label{prop:strong-module-monad-and-algebra-objects}
  Let \(\cat{C}\) be a monoidal category.
  There is a bijection between strong \(\cat{C}\)\hyp{}module monads on \(\cat{C}\)
  and algebra objects in \(\cat{C}\).
\end{proposition}
\begin{proof}
  By \cref{prop:bicategorical-yoneda-correspondence-monads-and-algebras},
  there is an equivalence of left module categories
  \begin{align*}
    \mathsf{Str}\cat{C}\mathsf{Mod}(\cat{C}, \cat{C}) \iso \cat{C}^{\tensorop}, \qquad
    T \mapsto T1, \qquad
    \blank \otimes A \longmapsfrom A,
  \end{align*}
  and by \cref{ex:monoid-iff-tensoring-is-monad}, \(A\) is an algebra object if and only if \(\blank \otimes A\) is a monad.

  In particular, if \((A, \tau, \upsilon)\) is an algebra object in \(\cat{C}\),
  then \((\blank \otimes A,\, \blank \otimes \tau,\, \blank \otimes \upsilon)\) becomes a strong \(\cat{C}\)\hyp{}module monad on \(\cat{C}\).
  The strong \(\cat{C}\)\hyp{}module structure is given by the associator of \(\cat{C}\):
  \[
    {(\blank \otimes A)}_{\mathsf{a};X} \defeq X \otimes (\blank \otimes A) \iso (X \otimes \blank) \otimes A, \qquad\text{for all \(X \in \cat{C}\)}.
  \]

  If \((T, \mu, \eta)\) is a strong \(\cat{C}\)\hyp{}module monad,
  then evaluating at the monoidal unit yields an algebra object in \(\cat{C}\),
  with multiplication \(\mu_1\) and unit \(\eta_1\).
\end{proof}

In particular, the category of right \(A\)\hyp{}modules and the Eilenberg–Moore category of \(\blank \otimes A\) coincide as \(\cat{C}\)\hyp{}module categories.

\begin{remark}\label{rem:comodule-monads-induce-module-categories}
  Let \(B \from \cat{C} \to \cat{C}\) be a bimonad and  \((C, C_{\mathsf{a}}) \from \cat{M} \to \cat{M}\) a comodule monad over it.
  The coaction of \(C\) allows us to define an action
  \[
    \ract \from \cat{M}^{C}\times \cat{C}^B \to \cat{M}^C.
  \]
  For any two modules \((m, \nabla_{m}) \in \cat{M}^C\) and \((x, \nabla_{x}) \in \cat{C}^B\), it is given by
  \[
    (m, \nabla_{m}) \ract (x, \nabla_{x}) \defeq \big(m \ract x, (\nabla_{m}\ract \nabla_{x}) \circ C_{\mathsf{a};m,x} \big).
  \]
  The axioms of the coaction of \(B\) on \(C\) translate precisely to the compatibility of the action of \(\cat{C}^B\) on \(\cat{M}^{C}\) with the tensor product and unit of \(\cat{C}^B\).
\end{remark}

\begin{definition}\label{def:comodule-adjunction}
  Let \(\stdadj\) be an oplax monoidal adjunction.
  Suppose that \(\adj{G}{V}{\cat{M}}{\cat{N}}\) is an adjunction
  such that \(G\) is an \(F\)-comodule functor,
  and \(V\) is a \(U\)-comodule functor.
  The pair \((G \dashv V,\ F \dashv U)\) is a \emph{comodule adjunction}
  if the following identities hold:
  \[
    \begin{mytikzcd}[ampersand replacement=\&]
	{m \ract x} \& {VG(m \ract x)} \& {GV(m \ract x)} \& {G(Vm \ract Ux)} \\
	{VGm \ract UFx} \& {V(Gm \ract Fx)} \& {m \ract x} \& {GVm \ract FUx}
	\arrow["{\eta^{(G \lad V)}_{m\ract x}}", from=1-1, to=1-2]
	\arrow["{VG_{\mathsf{a};m,x}}", from=1-2, to=2-2]
	\arrow["{V_{\mathsf{a};Gm,Fx}}", from=2-2, to=2-1]
	\arrow["{\eta^{(G \lad V)}_m \ract \eta^{(F \lad U)}_x}"', from=1-1, to=2-1]
	\arrow["{\varepsilon^{(G \lad V)}_{m \ract x}}"', from=1-3, to=2-3]
	\arrow["{GV_{\mathsf{a};m,x}}", from=1-3, to=1-4]
	\arrow["{G_{\mathsf{a};Vm,Ux}}", from=1-4, to=2-4]
	\arrow["{\varepsilon^{(G \lad V)}_m \ract \varepsilon^{(F \lad U)}_x}", from=2-4, to=2-3]
      \end{mytikzcd}
  \]
\end{definition}

\begin{example}\label{ex:module-adjunction-as-comodule-adjunction}
  If \(F = U = \Id_{\cat{M}}\) for a (left) \(\cat{C}\)\hyp{}module category \(\cat{M}\),
  then we say that \(G \lad V\) is an \emph{oplax \(\cat{C}\)\hyp{}module adjunction}.
  More explicitly, we have that \(G\) and \(V\) are oplax \(\cat{C}\)\hyp{}module functors
  such that the unit and counit of the adjunction are \(\cat{C}\)\hyp{}module transformations.

  Analogously to this, one can define the notion of a \emph{lax \(\cat{C}\)\hyp{}module adjunction}.
\end{example}

\Cref{fig:comodule-adjunction-unit-and-counit}
makes plain that the conditions required by \cref{def:comodule-adjunction}
are analogous to those stated in \cref{fig:comonoidal-adjunction-epsilon}.
\begin{figure}[htbp]
  \centering
  \tikzfig{comodule-adjunction-unit-and-counit}
  \scaption[.5cm]{%
    String diagrammatic conditions for a comodule adjunction.%
    \label{fig:comodule-adjunction-unit-and-counit}%
  }
\end{figure}

\begin{example}\label{ex:comodule-monad-over-bimonad}
  The philosophy that monads and adjunctions are two sides of the same coin
  extends to comodule functors.
  Suppose that we have an oplax monoidal adjunction \(\stdadj\)
  and over it a comodule adjunction \(\adj{G}{V}{\cat{M}}{\cat{N}}\).
  By~\cite[Proposition~4.3.1]{Aguiar2012},
  the bimonad \(B \defeq UF\) admits a coaction on the monad \(C \defeq VG\);
  for all \(m \in \cat{M}\) and \(x \in \cat{C}\) it is given by
  \[
    VG(m \ract x) \xrightarrow{VG_{\mathsf{a};m,x}} V(Gm\ract Fx) \xrightarrow{V_{\mathsf{a};Gm,Fx}} VGm \ract UFx = Cm\ract Bx.
  \]
\end{example}

\subsection{Reconstruction for comodule monads}

\textsc{The next theorem offers a way}
to reconstruct comodule monads from their module categories.
It extends~\cite[Proposition~4.1.2]{Aguiar2012}
and~\cite[Lemma~7.10]{Turaev2017}.
Recall \cref{def:lift-to-monoidal-adjunction}, and note that
one can analogously define lifts of adjunctions to comodule adjunctions.

\begin{theorem}\label{thm:comodule-adjunctions-correspond-to-strong-comodule-functors}
  Let \(\cat{C}\) and \(\cat{D}\) be monoidal categories,
  and suppose that \(\cat{M}\) and \(\cat{N}\) are right \(\cat{C}\)\hyp{} and \(\cat{D}\)\hyp{}module categories, respectively.
  Let \(\stdadj\) be an oplax monoidal adjunction.
  Lifts of an adjunction \(\adj{G}{V}{\cat{M}}{\cat{N}}\) to a comodule adjunction are in bijection with lifts of \(V \from \cat{N} \to \cat{M}\) to a strong \(U\)-comodule functor.
\end{theorem}
\begin{proof}
  Let \(G \dashv V\) be a comodule adjunction over \(F \dashv U\).
  Define the inverse \(V_{\mathsf{a}}^{-1}\from V\blank \ract U\bblank \nt V(\blank \ract \bblank)\) of \(V_{\mathsf{a}}\) by
  \begin{equation}\label[diagram]{eq:inverse-of-coaction}
    \tikzfig{inverse-of-coaction}
  \end{equation}
  Using that \(G\) and \(V\) are part of a comodule adjunction,
  a straightforward computation proves \(V_{\mathsf{a}}^{-1} \circ V_{\mathsf{a}} = \id_{V(\blank \,\ract\, \bblank)}\):
  \[
    \tikzfig{inverse-of-coaction-calc}
  \]
  A similar strategy can be used to show that \(V_{\mathsf{a}} \circ V_{\mathsf{a}}^{-1} = \id_{V\blank \,\ract\, U\bblank}\).
  Thus, \(V_{\mathsf{a}}\) is a natural isomorphism and therefore \(V\) is a strong comodule functor.

  Conversely, suppose that \(V \from \cat{N} \to \cat{M}\) is a strong comodule functor over \(U\).
  Define an arrow \(G_{\mathsf{a}}\from G(\blank \ract \bblank) \nt G(\blank) \ract F(\bblank)\) by
  \begin{equation}\label[diagram]{eq:coaction-via-counit-unit-and-inverse}
    \tikzfig{coaction-via-counit-unit-and-inverse}
  \end{equation}
  By~\cite[Lemma~7.10]{Turaev2017}, the comultiplication and counit of \(F\) are given by
  \[
    F(x \otimes y)
    \xrightarrow{F(\eta_x \otimes \eta_y)} F(UFx \otimes UFy)
    \xrightarrow{F U_{-2; Fx, Fy}} FU(Fx \otimes Fy)
    \xrightarrow{\varepsilon_{Fx \otimes Fy}} Fx \otimes Fy,
  \]
  \[
    F1 \xrightarrow{FU_{-0}} FU1 \xrightarrow{\ \varepsilon_1\ } 1,
    \qquad\text{for all \(x,y \in \cat{C}\)}.
  \]
  Note that, graphically, \(F_2\) looks just like \cref{eq:coaction-via-counit-unit-and-inverse},
  with black strings taking the place of blue ones.

  We will show that \(G\) is a comodule functor over \(F\).
  \Cref{fig:ga-is-coassociative} shows that \(G_{\mathsf{a}}\from G(\blank \ract \bblank) \nt G(\blank) \ract F(\bblank)\) is
  coassociative in the sense of \cref{fig:coassoc-and-counital-comodule-functor}.
  \begin{figure}[htbp]
    \centering
    \tikzfig{ga-is-coassociative}\vspace{-0.3cm}% XXX
    \scaption[-1cm]{%
      The coassociativity condition of \(G\).%
      \label{fig:ga-is-coassociative}%
    }
  \end{figure}

  The fact that \(G_{\mathsf{a}}\) is counital follows from
  \[
    \tikzfig{ga-is-counital}
  \]

  A straightforward computation proves that the unit of the adjunction \(G \dashv V\) satisfies
  the axioms displayed in \cref{fig:comodule-adjunction-unit-and-counit}:
  \[
    \tikzfig{unit-is-comodule-adjunction}
  \]
  A similar argument for the counit shows that \(G \dashv V\) is a comodule adjunction.

  To see that these constructions are inverse to each other,
  first suppose that we have a comodule adjunction \(G \adjoint V\).
  By utilising \(V_{\mathsf{a}}^{-1}\) as given in \cref{eq:inverse-of-coaction}, we obtain another coaction \(\lambda\) on \(G\),
  see \cref{eq:coaction-via-counit-unit-and-inverse}.
  Now, a direct computation shows that \(G_{\mathsf{a}} = \lambda\):
  \[
    \tikzfig{ga-is-lambda}
  \]

  The converse is clear:
  the map that associates a comodule adjunction \(G \lad V\) to any strong comodule structure on \(V\)
  preserves the coaction of \(V\).
\end{proof}

The proof of \cref{thm:comodule-adjunctions-correspond-to-strong-comodule-functors}
yields an analogue of \cref{prop:doctrinaladjunction},
which we formulate in the language of \(\cat{C}\)\hyp{}module functors.

\begin{porism}\label{porism:doctrinal-module-adjunctions}
  Let \(\cat{M}\) and \(\cat{N}\) be \(\cat{C}\)\hyp{} and \(\cat{D}\)\hyp{}module categories, respectively.
  For an adjunction \(\adj{F}{U}{\cat{M}}{\cat{N}}\),
  there is a bijection between oplax \(\cat{C}\)\hyp{}module structures on \(F\)
  and lax \(\cat{C}\)\hyp{}module structures on \(U\).
\end{porism}

\begin{corollary}\label{oplaxEMcorrespondence}
  Let \(\mathcal{M}\) be a \(\mathcal{C}\)\hyp{}module category,
  and let \(T\) be a monad on \(\mathcal{M}\).
  Then there exists a bijective correspondence\\
  \hspace{-2em}
  \scalebox{0.9}{\parbox{\linewidth}{%
      \vspace{-0.5em}
      \begin{align*}
        \big\{
        \text{oplax \(\mathcal{C}\)\hyp{}module monad structures on \(T\)}
        \big\}
        \;\leftrightarrow\,
        &\,\Bigg\{
          {\begin{aligned}
            &\text{\(\mathcal{C}\)\hyp{}module category structures on \(\cat{M}^T\)}\\
            &\text{such that \(U^T\) is a strict module functor}
          \end{aligned}}
          \Bigg\} \\
        \{\, T(\blank \triangleright \bblank) \xrightarrow{\;\;{}T_{\mathsf{a}}\;\;} \blank \triangleright T(\bblank) \,\}
        \mapsto &\,\{\, T(\blank \triangleright \bblank) \xrightarrow{\;\;{}T_{\mathsf{a}}\;\;} \blank \triangleright T(\bblank) \xrightarrow{\id \,\triangleright\, \mathrm{act}} \blank \triangleright \bblank \,\}\\
        \{\, U^T F^T(\blank \lact \bblank) \xrightarrow{U^T_{\mathsf{a}} \,\circ\, U^T F^T_{\mathsf{a}}} \blank \lact U^T F^T(\bblank) \,\}
        \longmapsfrom&\,\{\, U^T(\blank \lact \bblank) \xrightarrow{U^T_{\mathsf{a}}} \blank \lact U^T(\bblank) \,\}
      \end{align*}}}
\end{corollary}

\Cref{thm:comodule-adjunctions-correspond-to-strong-comodule-functors}
also yields a description of a comodule monad coaction in terms of its \EM{} adjunction;
\ie, a Tannaka--Krein reconstruction result in the spirit of \cref{prop: Moerdijk_reconstruction}.

\begin{theorem}\label{thm:comodule-monad-reconstruction}
  Let \(B\) be a bimonad on the monoidal category \(\cat{C}\),
  and \(T\) a monad on a right \(\cat{C}\)\hyp{}module category \(\cat{M}\).
  Coactions of \(B\) on \(T\) are
  in bijection with
  right actions of \(\cat{C}^B\) on \(\cat{M}^T\) such that \(U^{T}\) is a strict comodule functor over \(U^{B}\).
\end{theorem}
\begin{proof}
  Suppose \(\cat{C}^{B}\) acts from the right on \(\cat{M}^{T}\) such that \(U^{T}\) is a strict comodule functor.
  Due to \cref{thm:comodule-adjunctions-correspond-to-strong-comodule-functors},
  \(T = U^T F^T\) is a comodule monad via the coaction
  \[
    T_{\mathsf{a}}\from T(\blank \ract \bblank)
    \xrightarrow{U^T F^T_{\mathsf{a}}} U^T(F^T(\blank) \ract F^B(\bblank))
    \xrightarrow{\ U^T_{\mathsf{a}}\ } T(\blank) \ract B(\bblank),
  \]
  which is equal to \(U^T F^T_{\mathsf{a}}\), as \(U^T\) is a strict \(U^B\)-comodule functor.

  Conversely, if \(T\) is a comodule monad,
  then \(\cat{M}^{T}\) becomes a right \(\cat{C}^{B}\)-module category,
  with action given as in \cref{rem:comodule-monads-induce-module-categories}.

  As \(T_{\mathsf{a}}\) and the action of \(\cat{C}^B\) on \(\cat{M}^T\)
  determine the coaction \(F^T_{\mathsf{a}}\) of \(F^{T}\) uniquely,
  the two constructions are inverse to each other by
  \cref{thm:comodule-adjunctions-correspond-to-strong-comodule-functors}.
\end{proof}

The following result is dual to \cref{thm:comodule-adjunctions-correspond-to-strong-comodule-functors}
for the special case of \(\cat{C}\)\hyp{}module adjunctions,
where we study strong \(\mathcal{C}\)\hyp{}module structures on the left rather than the right adjoint.

\begin{proposition}\label{prop:dual-module-adjunctions}
  Let \(\cat{C}\) be a monoidal category,
  and suppose that \(\cat{M}\) and \(\cat{N}\) are left
  \(\cat{C}\)\hyp{}module categories.
  Given an adjunction \(\adj{F}{U}{\cat{M}}{\cat{N}}\),
  there is a one-to-one correspondence between
  lifts of \(F\) to a strong \(\cat{C}\)\hyp{}module functor,
  and lifts of \(F \dashv U\) to a lax \(\cat{C}\)\hyp{}module adjunction.
\end{proposition}
\begin{proof}
  Suppose that \(\adj{F}{U}{\cat{M}}{\cat{N}}\) is a lax \(\cat{C}\)\hyp{}module adjunction.
  Define the inverse of the action \(F_{\mathrm{a}} \from \blank \lact F(\bblank) \nt F(\blank \lact \bblank)\)
  by
  \[
    F(\blank \lact \bblank)
    \xrightarrow{F(\eta \,\lact\, \bblank)} F(UF(\blank) \lact \bblank)
    \xrightarrow{F U_{\mathrm{a}}} FU(F(\blank) \lact \bblank)
    \xrightarrow{\;\;\varepsilon\;\;} F(\blank) \lact \bblank.
  \]
  Conversely, given an adjunction \(\adj{F}{U}{\cat{M}}{\cat{N}}\)
  such that \(F\) is a strong \(\cat{C}\)\hyp{}module functor,
  define
  \[
    U(\blank) \lact \bblank
    \xrightarrow{\;\;\eta\;\;} UF(U(\blank) \lact \bblank)
    \xrightarrow{U F_{\mathrm{a}}^{-1}} U(FU(\blank) \lact \bblank)
    \xrightarrow{U(\varepsilon \,\lact\, \bblank)} U(\blank \lact \bblank).
  \]

  Reading the string diagrams in the proof of \cref{thm:comodule-adjunctions-correspond-to-strong-comodule-functors} upside down
  verifies the necessary coherence conditions,
  and that these two constructions are inverses of each other.
\end{proof}

We obtain a version of \cref{oplaxEMcorrespondence}
for the Kleisli category of a monad.

\begin{corollary}\label{laxKlCorrespondence}
  Let \(\mathcal{M}\) be a left \(\mathcal{C}\)\hyp{}module category,
  and suppose \(T\) to be a monad on \(\mathcal{M}\).
  There is a bijection between
  lax \(\cat{C}\)\hyp{}module monad structures on \(T\),
  and
  \(\cat{C}\)\hyp{}module category structures on \(\cat{M}_T\),
  such that \(F_T\) is a strict \(\cat{C}\)\hyp{}module functor.
\end{corollary}
\begin{proof}
  Let \(\cat{M}_T\) be a \(\cat{C}\)\hyp{}module category
  such that \(F_T\) is a strict \(\cat{C}\)\hyp{}module functor.
  Then
  \(T \defeq U_T F_T \from \cat{M} \to \cat{M}\)
  is a lax \(\cat{C}\)\hyp{}module monad by \cref{prop:dual-module-adjunctions}.

  Conversely, if \(T\) is a lax \(\cat{C}\)\hyp{}module monad,
  \(\cat{M}_T\) becomes a \(\cat{C}\)\hyp{}module category as follows:
  for \(x \in \cat{C}\) and \(m,n \in \cat{M}\),
  set \(x \lact_{\cat{M}_T} m \defeq x \lact_{\mathcal{M}} m\) on objects,
  and for a morphism \(f \from m \to Tn\) in \(\cat{M}_T\),
  we define the action of \(x \in \cat{C}\) by
  \[
    x \lact_{\cat{M}_T} f \defeq x \lact m \xrightarrow{x \,\lact\, f} x \lact Tn \xrightarrow{T_{\mathsf{a};x,n}} T(x \lact n).
  \]
  It is easy to check that these assignments define a \(\cat{C}\)\hyp{}module structure,
  for which \(F_T\) is a strict \(\cat{C}\)\hyp{}module functor.

  The \(\cat{C}\)\hyp{}module structure of \(\cat{M}_T\)
  and the lax module monad structure on \(T\) uniquely determine the \(\cat{C}\)\hyp{}module structure of \(F_T\).
  Hence, \(F_{T;\mathsf{a}}^{-1}\) is given as in \cref{prop:dual-module-adjunctions}
  and the two constructions are inverses of each other.
\end{proof}

Using these insights, we may now clarify the structure of comparison functors associated to comodule adjunctions.

\begin{proposition}\label{prop:comparison-functor-comodule-monad-comodule}
  Consider a comodule adjunction \(\adj{G}{V}{\cat{M}}{\cat{N}}\) over an oplax monoidal adjunction \(\stdadj\),
  and denote the associated comodule monad and bimonad by
  \(C \defeq VG \from \cat{M} \to \cat{M}\)
  and \(B \defeq UF \from \cat{C} \to \cat{C}\), respectively.

  Then the comparison functor \(K^{C} \from \cat{N} \to \cat{M}^C\) is
  a strong comodule functor over \(K^{B} \from \cat{D} \to \cat{C}^B\),
  and the following identities of comodule functors hold:
  \[
    U^C K^{C} = V \qquad\text{and}\qquad K^{C} G = F^C.
  \]
\end{proposition}
\begin{proof}
  We proceed analogously to~\cite[Theorem~2.6]{Bruguieres2007}.
  For any \(n \in \cat{N}\) we have \(K^{C}n= (Vn, V\varepsilon_{n})\) and a direct computation shows that the coaction of \(V\) lifts to a coaction of \(K^{C}\).
  That is,
  \[
    U^{C} K^C_{\mathsf{a}; n, y} = V_{\mathsf{a}; n, y}, \qquad \qquad \text{ for all } n \in \cat{N} \text{ and } y \in \cat{D}.
  \]
  Using that \(U^{C} \from \cat{M}^{C} \to \cat{M}\) is a faithful and conservative functor,
  one observes that \(K^{C}\) becomes a strong comodule functor in this manner.
  Furthermore, as \(U^{C}\) is a strict comodule functor,
  the coactions of \(U^{C}K^{C}\) and \(V\) coincide.
  Lastly, we compute the following, for any \(x \in \cat{C}\) and
  \(m \in \cat{M}\):
  \[\mathscale{0.98}{
      {(K^{C} G)}_{\mathsf{a};m, x}
      = {(U^C K^{C} G)}_{\mathsf{a};m,x}
      = {(VG)}_{\mathsf{a};m,x}
      = C_{\mathsf{a};m,x}
      = {(U^C F^C)}_{\mathsf{a};m,x}
      = F^C_{\mathsf{a};m,x}.
    }\qedhere\]
\end{proof}

\begin{proposition}\label{strongcomparisons}
  Let \(\cat{M}\) and \(\cat{N}\) be left \(\cat{C}\)\hyp{}module categories.
  \begin{itemize}
    \item Let \(\adj{F}{U}{\cat{M}}{\cat{N}}\) be an oplax \(\cat{C}\)\hyp{}module adjunction
    and consider the corresponding \(\cat{C}\)\hyp{}module category structure on \(\cat{M}^T\) of \cref{oplaxEMcorrespondence}.
    Then for \(T \defeq UF\) the comparison functor \(K^T\) is a strong \(\cat{C}\)\hyp{}module functor.
    \item Let \(\adj{F}{U}{\cat{M}}{\cat{N}}\) be
    a lax \(\cat{C}\)\hyp{}module adjunction
    and consider the corresponding \(\cat{C}\)\hyp{}module category structure on \(\cat{M}_T\) of \cref{laxKlCorrespondence}.
    Then for \(T \defeq UF\) the comparison functor \(K_T\) is a strong \(\cat{C}\)\hyp{}module functor.
  \end{itemize}
\end{proposition}

\begin{proposition}\label{EMMonadCorrespondence}
  Let \(\cat{M}\) be a \(\cat{C}\)\hyp{}module category,
  \(T\from \cat{M} \to \cat{M}\) an oplax \(\cat{C}\)\hyp{}module monad,
  and \(S\from \cat{M} \to \cat{M}\) a right adjoint to \(T\).
  Then the isomorphism \(\cat{M}^T \cong \cat{M}^S\) of Proposition~\ref{pyramidscheme} is a strict \(\cat{C}\)\hyp{}module isomorphism,
  where \(\cat{M}^T\) is endowed with the \(\cat{C}\)\hyp{}module structure of Proposition~\ref{oplaxEMcorrespondence},
  and similarly for \(\cat{M}^S\).
\end{proposition}
\begin{proof}
  Note that,
  by \cref{thm:comodule-adjunctions-correspond-to-strong-comodule-functors}
  the comonad \(S\from \cat{M}\to \cat{M}\) comes equipped with the following lax \(\cat{C}\)\hyp{}module structure,
  for all \(x \in \cat{C}\) and \(m \in \cat{M}\):
  \[
    S_{\mathsf{a}; x, m} \from
    x \lact S m
    \xrightarrow{\eta_{x \lact S m}} S T(x \lact S m)
    \xrightarrow{S T_{\mathsf{a}; x, S m}} S(x \lact TS m)
    \xrightarrow{S(x \lact \varepsilon_m)} S(x \lact m).
  \]
  Further, the isomorphism \(L\from \cat{M}^T \iso \cat{M}^S\) of \cref{pyramidscheme} is given by
  \[
    \big(m,\ \nabla_m\from Tm \to m\big) \mapsto \big(m,\ m \xrightarrow{\ \eta_m\ } S Tm \xrightarrow{\,S\nabla_m\,} S m\big).
  \]

  We have to prove that \(L(x \lact (m, \nabla_m)) = x \lact L((m, \nabla_m))\),
  for all \(x \in \cat{C}\) and \((m, \nabla_m) \in \cat{M}^T\).
  This is equivalent to the equality of
  \[
    \big(
    x \lact m,\
    x \lact m
    \xrightarrow{\eta_{x \lact m}} S T(x \lact m)
    \xrightarrow{S T_{\mathsf{a};x, m}} S(x \lact Tm)
    \xrightarrow{S(x \lact \nabla_m)} S(x \lact m)
    \big)
  \]
  and
  \begin{align*}
    \big(
    x \lact m,\
    x \lact m
    &\xrightarrow{x \lact \eta_m} x \lact S Tm
      \xrightarrow{x \lact S \nabla_m} x \lact S m
      \xrightarrow{\eta_{x \lact S m}} S T(x \lact S m) \\
    &\xrightarrow{S{(T_{\mathsf{a}})}_{x,S(m)}} S(x \lact TS m)
      \xrightarrow{S(x \lact \varepsilon_m)} S(x \lact m)
      \big).
  \end{align*}
  This is evidenced by the following commutative diagram:
  \[
    \begin{mytikzcd}[ampersand replacement=\&]
      {x \lact m} \& {S T(x\lact m)} \&\& {S(x\lact Tm)} \\
      {x\lact S Tm} \& {S T(x \lact S Tm)} \& {S(x \lact TS T m)} \& {S(x\lact Tm)} \\
      {x\lact S m} \& {S T(x \lact S m)} \& {S(x\lact TS m)} \& {S(x \lact m)}
      \arrow[""{name=0, anchor=center, inner sep=0}, "{{{\eta_{x\lact m}}}}", from=1-1, to=1-2]
      \arrow["{{{x\lact \eta_m}}}"', from=1-1, to=2-1]
      \arrow[""{name=1, anchor=center, inner sep=0}, "{{{S T_{\mathsf{a};x,m}}}}", from=1-2, to=1-4]
      \arrow["{{{S T(x \lact \eta_m)}}}", from=1-2, to=2-2]
      \arrow[""{name=2, anchor=center, inner sep=0}, "{{{S(x \lact T\eta_m)}}}"', from=1-4, to=2-3]
      \arrow[equals, from=1-4, to=2-4]
      \arrow[""{name=3, anchor=center, inner sep=0}, "{{{\eta_{x \lact S Tm}}}}"', from=2-1, to=2-2]
      \arrow["{{{x \lact S \nabla_m}}}"', from=2-1, to=3-1]
      \arrow[""{name=4, anchor=center, inner sep=0}, "{{{S T_{\mathsf{a}; x, S T m}}}}", from=2-2, to=2-3]
      \arrow["{{{S T(x \lact S \nabla_m)}}}", from=2-2, to=3-2]
      \arrow[""{name=5, anchor=center, inner sep=0}, "{{{S(x \lact \varepsilon_{Tm})}}}"', from=2-3, to=2-4]
      \arrow["{{{S(x \lact TS \nabla_m)}}}", from=2-3, to=3-3]
      \arrow["{{{S(x \lact \nabla_m)}}}", from=2-4, to=3-4]
      \arrow[""{name=6, anchor=center, inner sep=0}, "{{{\eta_{x \lact S m}}}}"', from=3-1, to=3-2]
      \arrow["{{{S T_{\mathsf{a};x,S m}}}}"', from=3-2, to=3-3]
      \arrow["{{{S(x \lact \varepsilon_{m})}}}"', from=3-3, to=3-4]
      \arrow["{\mathsf{nat}\ \eta}"{description}, draw=none, from=0, to=3]
      \arrow["{\mathsf{nat}\ T_{\mathsf{a}}}"{description}, draw=none, from=1, to=4]
      \arrow["{{\mathsf{adj}}}"{description}, draw=none, from=2, to=2-4]
      \arrow["{\mathsf{nat}\ \eta}"{description}, draw=none, from=3, to=6]
      \arrow["{\mathsf{nat}\ T_{\mathsf{a}}}"{description, pos=0.7}, draw=none, from=4, to=3-3]
      \arrow["{\mathsf{nat}\ \varepsilon}"{description, pos=0.7}, draw=none, from=5, to=3-4]
    \end{mytikzcd}
  \]
\end{proof}

In an analogous way to \cref{EMMonadCorrespondence},
one obtains the following result.

\begin{proposition}\label{KleisliMonadCorrespondence}
  Let \(\cat{M}\) be a \(\cat{C}\)\hyp{}module category,
  \(T\from \cat{M}\to \cat{M}\) a lax \(\cat{C}\)\hyp{}module monad,
  and \(L\from \cat{M} \to \cat{M}\) a left adjoint to \(T\).
  Then the isomorphism \(\cat{M}_T \cong \cat{M}_{L}\) of \cref{einerKleinerKleisliÄquivalenz}
  is a \(\cat{C}\)\hyp{}module isomorphism,
  where \(\cat{M}_T\) and \(\cat{M}_L\) are
  endowed with the \(\cat{C}\)\hyp{}module structure of \cref{laxKlCorrespondence}.
\end{proposition}

%%% Local Variables:
%%% mode: LaTeX
%%% TeX-master: "../main"
%%% TeX-command-extra-options: "-shell-escape -fmt=prec.fmt -file-line-error -synctex=1"
%%% End:
